%% ICEBA2026 — Springer Nature Proceedings (sn-jnl). BAN TIENG VIET (de doc/duyet).
%% Ban tieng Anh (nop chinh thuc) o file iceba_astro_en.tex cung thu muc.
%%
%% >>> Bien dich bang pdfLaTeX. Tieng Viet co dau nho babel + [T5]{fontenc}.
%%     Neu dau bi loi -> doi Compiler sang XeLaTeX va thay fontenc/babel bang:
%%       \usepackage{fontspec}
%%       \usepackage{polyglossia}\setmainlanguage{vietnamese}
\documentclass[pdflatex,sn-mathphys-num]{sn-jnl}

\usepackage{lmodern}
\usepackage[utf8]{inputenc}
\usepackage[T5]{fontenc}
\usepackage[vietnamese]{babel}

\addto\captionsvietnamese{%
  \renewcommand{\abstractname}{Abstract}%
  \renewcommand{\tablename}{Table}%
  \renewcommand{\figurename}{Fig.}%
  \renewcommand{\refname}{References}%
  \renewcommand{\bibname}{References}%
}

\usepackage{graphicx}
\graphicspath{{figures/}{./}{../figures/}}
\usepackage{multirow}
\usepackage{amsmath,amssymb,amsfonts}
\usepackage{booktabs}

\raggedbottom

\begin{document}

\title[Hiệu chỉnh conformal thích nghi cho đếm nguồn thiên văn]{Hiệu chỉnh conformal trực tuyến thích nghi
cho đếm nguồn thiên văn dưới dịch chuyển độ sâu khảo sát}

\author[1]{\fnm{Thi Thu Hiep} \sur{Dinh}}\email{dinhthithuhiep1@gmail.com}
\author*[1]{\fnm{Viet Hang} \sur{Duong}}\email{hangdv@uit.edu.vn}

\affil*[1]{\orgdiv{Khoa Khoa học Máy tính}, \orgname{Trường Đại học Công nghệ Thông tin (UIT), ĐHQG-HCM},
\orgaddress{\city{TP. Hồ Chí Minh}, \country{Việt Nam}}}

\abstract{%
Đếm nguồn thiên văn là một tác vụ nền tảng trong thiên văn quan trắc, làm cơ sở cho các phép đo mật độ nguồn,
thống kê quần thể và cấu trúc quy mô lớn. Tuy nhiên, hầu hết các pipeline phát hiện nguồn chỉ đưa ra ước lượng
điểm mà không định lượng độ bất định, một hạn chế trở nên nghiêm trọng khi dữ liệu triển khai khác với dữ liệu
hiệu chỉnh do thay đổi độ sâu khảo sát hoặc điều kiện quan trắc. Conformal prediction cung cấp khoảng dự đoán
với bảo đảm coverage hữu hạn mẫu dưới giả định trao đổi được, song hiệu chỉnh tĩnh suy giảm dưới dịch chuyển
phân phối, trong khi bài toán đếm đa lớp còn đòi hỏi phủ đồng thời toàn bộ vector số đếm. Bài báo áp dụng
Adaptive PB-JCI Online, một khung hiệu chỉnh conformal trực tuyến vốn được phát triển cho bài toán đếm tế bào,
sang bài toán đếm nguồn thiên văn. Phương pháp kết hợp thang phương sai Poisson–Binomial, hiệu chỉnh phủ đồng
thời qua một điểm nonconformity dạng thống kê cực đại, và một cơ chế cập nhật cửa sổ trượt thích nghi. Các thí
nghiệm được tiến hành trên $61\,423$ nguồn thật từ Sloan Digital Sky Survey DR17, sử dụng danh mục quang trắc
chính thức làm backbone xác suất và một dịch chuyển độ sâu khảo sát có kiểm soát được tạo ra bằng suy giảm tỉ
số tín hiệu trên nhiễu. Phương pháp đề xuất duy trì joint coverage hợp lệ khoảng $91\%$ qua mọi mức dịch
chuyển, trong khi hiệu chỉnh tĩnh suy giảm từ $85\%$ xuống $9\%$. So với một số baseline conformal trực tuyến
tiêu biểu, Adaptive PB-JCI đạt coverage hợp lệ đồng thời có điểm interval Winkler thấp nhất trong các phương
pháp được đánh giá. Các kết quả này cho thấy một khung hiệu chỉnh đếm được phát triển cho ứng dụng khoa học
sự sống có thể chuyển giao hiệu quả sang bài toán đếm nguồn thiên văn dưới dịch chuyển phân phối.
}

\keywords{conformal prediction, uncertainty quantification, count calibration, distribution shift,
astronomical source counting}

\maketitle

%%===========================================================%%
\section{Giới thiệu}\label{sec:intro}

Đếm và phân loại nguồn thiên văn thành các lớp như sao và thiên hà là một tác vụ nền tảng trong vật lý thiên
văn quan trắc. Nhiều đại lượng khoa học quan trọng, chẳng hạn mật độ nguồn theo vùng trời, hàm phân bố quần
thể và tỉ lệ giữa các lớp nguồn, được suy ra trực tiếp từ số lượng và thành phần của các nguồn được phát
hiện. Các khảo sát quang học hiện đại tạo ra danh mục gồm hàng triệu nguồn cùng với xác suất phân loại và các
đặc trưng quang trắc đo được, từ đó số đếm theo từng lớp nguồn được sử dụng trong nhiều phân tích thống kê
thiên văn.

Tuy nhiên, các hệ đếm dựa trên phát hiện thường chỉ cung cấp một ước lượng điểm mà không định lượng độ bất
định của kết quả. Trong thực tế, điều quan trọng không chỉ là số lượng nguồn được dự đoán mà còn là mức độ
tin cậy đi kèm với dự đoán đó. Nhu cầu này trở nên đặc biệt cấp thiết khi dữ liệu triển khai khác với dữ liệu
hiệu chỉnh do thay đổi thiết bị, độ sâu khảo sát hoặc điều kiện quan trắc. Trong những trường hợp như vậy,
một mô hình vẫn có thể đạt sai số trung bình chấp nhận được nhưng tạo ra các dự đoán bị hiệu chỉnh kém trên
dữ liệu ngoài miền.

Conformal prediction là một khung thống kê cung cấp khoảng dự đoán với bảo đảm độ phủ hữu hạn mẫu dưới giả
định trao đổi được giữa dữ liệu hiệu chỉnh và dữ liệu kiểm tra. Tuy nhiên, việc áp dụng conformal prediction
cho bài toán đếm đa lớp gặp phải hai thách thức chính. Thứ nhất, các phương pháp hiệu chỉnh tĩnh dựa trên một
tập hiệu chỉnh cố định thường suy giảm độ phủ khi dữ liệu triển khai xuất hiện dịch chuyển phân phối. Thứ
hai, mục tiêu của bài toán đếm đa lớp không chỉ là bảo đảm độ phủ cho từng lớp riêng lẻ mà còn phải bảo đảm
độ phủ đồng thời cho toàn bộ vector số đếm.

Adaptive PB-JCI Online là một phương pháp hiệu chỉnh conformal trực tuyến được phát triển gần đây cho các bài
toán đếm dựa trên xác suất~\cite{dinh2026adaptive}. Nghiên cứu này tập trung đánh giá tính tổng quát và khả
năng chuyển giao của phương pháp sang một miền ứng dụng khoa học khác biệt đáng kể, cụ thể là bài toán đếm
nguồn thiên văn.

Các đóng góp chính của bài báo bao gồm:
\begin{itemize}
\item Chuyển giao Adaptive PB-JCI Online từ bài toán đếm tế bào sang bài toán đếm nguồn thiên văn trên một
tập dữ liệu và backbone phân loại hoàn toàn khác.
\item Xây dựng một benchmark dựa trên dữ liệu thực từ Sloan Digital Sky Survey (SDSS) DR17 với các mức dịch
chuyển độ sâu khảo sát được kiểm soát cùng một tập baseline conformal trực tuyến đại diện.
\item Chứng minh thực nghiệm rằng Adaptive PB-JCI Online là phương pháp duy nhất duy trì joint coverage hợp
lệ trên toàn bộ dải dịch chuyển khảo sát, đồng thời đạt điểm Winkler thấp nhất trong số các phương pháp được
đánh giá.
\end{itemize}

Conformal prediction cung cấp khoảng dự đoán không phụ thuộc phân phối với bảo đảm độ phủ hữu hạn mẫu dưới
giả định trao đổi được giữa dữ liệu hiệu chỉnh và dữ liệu kiểm tra~\cite{vovk2005algorithmic}. Trong bối
cảnh dự đoán trực tuyến dưới dịch chuyển phân phối, nhiều phương pháp đã được đề xuất nhằm duy trì coverage
theo thời gian, bao gồm Adaptive Conformal Inference (ACI)~\cite{gibbs2021adaptive},
FACI~\cite{gibbs2024conformal}, NexCP~\cite{barber2023conformal}, SAOCP~\cite{bhatnagar2023improved},
COP~\cite{hu2026conformal} và các biến thể rolling-origin conformal prediction~\cite{halkiewicz2026rolling}.
Tuy nhiên, các phương pháp này chủ yếu hướng tới dự đoán đơn biến, trong khi bài toán đếm đa lớp đòi hỏi độ
phủ đồng thời cho toàn bộ vector số đếm. Khoảng trống này là động lực cho nghiên cứu hiện tại.

%%===========================================================%%
\section{Phương pháp}\label{sec:method}

Adaptive PB-JCI Online là một tầng hiệu chỉnh conformal được xây dựng trên đầu ra của một backbone phát hiện
và phân loại nguồn. Phương pháp gồm ba thành phần chính: (i) một thang phương sai dựa trên phân phối
Poisson--Binomial để mô hình hóa độ bất định của số đếm, (ii) một điểm nonconformity dựa trên thống kê cực
đại nhằm bảo đảm độ phủ đồng thời cho toàn bộ vector số đếm, và (iii) một cơ chế hiệu chỉnh trực tuyến sử
dụng cửa sổ trượt thích nghi để duy trì coverage dưới dịch chuyển phân phối. Phần này trình bày ngắn gọn các
thành phần chính của phương pháp nhằm bảo đảm tính tự chứa của bài báo.

\subsection{Số đếm mềm và thang phương sai Poisson--Binomial}
Với mỗi vùng quan sát (\emph{field}), backbone cung cấp xác suất phát hiện $s_i$ cho từng thực thể ứng viên
$i$, cùng vector xác suất lớp $p_i[k]$. Thay vì sử dụng số đếm cứng, số đếm dự đoán của lớp $k$ được xác định
bằng tổng kỳ vọng
\[
n_{\mathrm{pred}}[k]=\sum_i s_i\,p_i[k].
\]
Độ lệch chuẩn tương ứng được tính theo
\begin{equation}
\sigma[k]=\sqrt{\sum_i \big(s_i p_i[k]\big)\big(1-s_i p_i[k]\big)}. \label{eq:soft}
\end{equation}
Đại lượng $\sigma[k]$ tương ứng với độ lệch chuẩn của tổng các biến Bernoulli độc lập, tức phân phối
Poisson--Binomial. Nhờ đó, độ rộng khoảng dự đoán được điều chỉnh theo mức độ bất định cục bộ của từng vùng
quan sát: các vùng chứa nhiều thực thể có xác suất không chắc chắn sẽ tạo ra giá trị $\sigma[k]$ lớn hơn,
trong khi các vùng có dự đoán rõ ràng sẽ có giá trị nhỏ hơn.

\subsection{Nonconformity cho độ phủ đồng thời}
Mục tiêu của bài toán là bảo đảm độ phủ đồng thời cho toàn bộ vector số đếm thay vì cho từng lớp riêng lẻ. Do
đó, điểm nonconformity được định nghĩa dưới dạng thống kê cực đại chuẩn hóa qua tất cả các lớp:
\begin{equation}
S_t=\max_k \frac{\big|\,g_t[k]-n_{\mathrm{pred}}[k]\,\big|}{\sigma[k]}, \label{eq:maxstat}
\end{equation}
trong đó $g_t[k]$ là số đếm thực của lớp $k$. Cho $q$ là phân vị bậc $(1-\alpha)$ của các điểm nonconformity
trên tập hiệu chỉnh, khoảng dự đoán cho từng lớp được xác định bởi
\[
\big[\,\max(0,\ n_{\mathrm{pred}}[k]-q\,\sigma[k]),\ \ n_{\mathrm{pred}}[k]+q\,\sigma[k]\,\big].
\]
Việc sử dụng thống kê cực đại cho phép chuyển yêu cầu độ phủ đồng thời của toàn bộ vector số đếm thành một
điều kiện hiệu chỉnh duy nhất.

\subsection{Cửa sổ hiệu chỉnh thích nghi trực tuyến}
Để thích ứng với dữ liệu thay đổi theo thời gian, ngưỡng hiệu chỉnh $q$ được cập nhật trực tuyến thông qua
một cửa sổ hiệu chỉnh có kích thước biến thiên. Gọi $w_t$ là kích thước cửa sổ tại thời điểm $t$, và
$\hat c_t$ là độ phủ quan sát được trên một cửa sổ giám sát gần nhất. Kích thước cửa sổ được cập nhật theo
quy tắc
\begin{equation}
w_{t+1}=
\begin{cases}
\max\!\big(w_{\min},\ \lfloor \rho^{-}\,w_t \rfloor\big), & \hat c_t < 1-\alpha,\\[4pt]
\min\!\big(w_{\max},\ \lfloor \rho^{+}\,w_t \rfloor\big), & \hat c_t > 1-\alpha+\beta,\\[4pt]
w_t, & \text{ngược lại},
\end{cases}
\label{eq:window}
\end{equation}
trong đó $\rho^{-}<1<\rho^{+}$ và $\beta>0$ là một biên trễ. Sau khi cập nhật $w_{t+1}$, ngưỡng $q$ được tính
lại từ phân vị bậc $(1-\alpha)$ của $w_{t+1}$ điểm nonconformity gần nhất. Khi coverage giảm dưới mức mục
tiêu, cửa sổ được thu hẹp để ngưỡng hiệu chỉnh phản ứng nhanh hơn với dữ liệu mới. Ngược lại, khi coverage
vượt mục tiêu một lượng đủ lớn, cửa sổ được mở rộng nhằm giảm dao động của ước lượng. Cơ chế cập nhật trực
tuyến này là thành phần then chốt giúp duy trì độ phủ dưới dịch chuyển phân phối, trong khi thang phương sai
Poisson--Binomial chủ yếu đóng vai trò chuẩn hóa độ rộng khoảng dự đoán giữa các vùng quan sát có mức độ bất
định khác nhau.

%%===========================================================%%
\section{Dữ liệu và thiết lập thực nghiệm}\label{sec:data}

\subsection{Dữ liệu SDSS DR17}
Nghiên cứu sử dụng dữ liệu từ Sloan Digital Sky Survey Data Release~17 (SDSS
DR17)~\cite{sdss_dr17,sdss_technical}, một khảo sát quang học trong năm băng $u,g,r,i,z$. Sau khi lọc chất
lượng (\texttt{mode=1}, \texttt{clean=1}), giữ hai lớp sao/thiên hà và chỉ giữ các nguồn có
\texttt{psfFlux\_r} và \texttt{psfFluxIvar\_r} dương (để tính được SNR), tập dữ liệu thu được gồm $61\,423$
nguồn ($25\,270$ sao và $36\,153$ thiên hà) với median tỉ số tín hiệu trên nhiễu (SNR) khoảng $11.1$. Vùng trời được chia thành $841$ ô lưới đóng vai
trò các đơn vị đếm độc lập (field); các ô chỉ số chẵn tạo tập hiệu chỉnh (CAL, $421$ field) và các ô chỉ số
lẻ tạo tập kiểm tra (TEST, $420$ field), bảo đảm sự tách biệt không gian giữa hai tập.

\subsection{Backbone quang trắc}
Mục tiêu của nghiên cứu là đánh giá hiệu chỉnh số đếm thay vì phát hiện nguồn, do đó pipeline quang trắc SDSS
được sử dụng như một backbone cố định, không huấn luyện lại. Từ danh mục, mỗi nguồn cung cấp các đại lượng
thật gồm flux, tỉ số tín hiệu trên nhiễu $\mathrm{SNR}_i=\texttt{psfFlux}_i\sqrt{\texttt{psfFluxIvar}_i}$ và
nhãn phân loại sao/thiên hà. Cần lưu ý rằng trong SDSS, \texttt{probPSF} là một \emph{cờ nhị phân} ($0$ hoặc
$1$) trùng với nhãn \texttt{type}, chứ không phải một xác suất được hiệu chỉnh; do đó việc gán lớp ở chế độ
khảo sát sâu là gán cứng.

Độ bất định theo từng nguồn, tức thành phần sinh ra thang phương sai Poisson--Binomial, được mô hình hóa trên
SNR thật theo hai cơ chế. Thứ nhất, do danh mục chỉ chứa các nguồn đã phát hiện nên xác suất phát hiện không
có sẵn và được lấy bằng hàm completeness logistic $s_i=\big[1+\exp(-(\mathrm{SNR}_i-\tau)/\kappa)\big]^{-1}$
với $\tau=5$ (ngưỡng $5\sigma$) và $\kappa=1.5$. Thứ hai, dưới dịch chuyển độ sâu (mô tả bên dưới), nhãn lớp
nhị phân được làm mờ dần về $0.5$ khi SNR giảm, phản ánh việc phân loại trở nên kém tin cậy đối với các nguồn
mờ. Xác suất phát hiện và xác suất lớp thu được được dùng để tính số đếm mềm và thang phương sai
trong~\eqref{eq:soft}; ở chế độ sâu, thang phương sai này chủ yếu phản ánh bất định phát hiện. Như vậy đây là
một thiết lập \emph{bán tổng hợp}: các đại lượng quang trắc là thật, còn xác suất phát hiện và độ mờ phân
loại là thành phần được mô hình hóa, neo trên SNR thật.

\subsection{Dịch chuyển độ sâu khảo sát}
Để tạo dịch chuyển phân phối có kiểm soát, SNR của từng nguồn được suy giảm theo
\[
\mathrm{SNR}_{\mathrm{shallow}}=\mathrm{SNR}_{\mathrm{deep}}/\texttt{DEPTH}.
\]
Khi \texttt{DEPTH} tăng, nhiều nguồn tiến gần hoặc rơi xuống dưới ngưỡng phát hiện, làm tăng độ bất định phân
loại và suy giảm hiệu năng đếm. Tập hiệu chỉnh ứng với chế độ khảo sát sâu và tập kiểm tra ứng với chế độ
nông; \texttt{DEPTH} vì vậy đóng vai trò điều khiển cường độ dịch chuyển trong toàn bộ thí nghiệm.

\subsection{Thiết lập và chỉ số đánh giá}
Ngưỡng completeness cho các thí nghiệm chính được chọn là ngưỡng tạo ra dịch chuyển đủ mạnh để hiệu chỉnh
tĩnh suy giảm rõ trên tập kiểm tra trong khi vẫn giữ coverage lành mạnh trên tập hiệu chỉnh. Adaptive PB-JCI
được so sánh với Static CP, ACI, NexCP, FACI, SAOCP, COP, Rolling-Origin CP và PB-JCI Fixed; tất cả được
triển khai trên cùng một thống kê nonconformity dạng cực đại nhằm đánh giá joint coverage cho $K{=}2$ lớp.
Các chỉ số báo cáo gồm joint coverage, coverage riêng từng lớp, độ rộng trung bình của khoảng và điểm
Winkler~\cite{winkler1972decision} (càng thấp càng tốt). Coverage mục tiêu là $90\%$ ($\alpha=0.1$); các tham
số cửa sổ thích nghi được giữ cố định trong mọi thí nghiệm ($w_{\min}=40$, $w_{\max}=300$, $\rho^{-}=0.9$,
$\rho^{+}=1.05$, $\beta=0.03$, cửa sổ giám sát $50$ bước). Mỗi cấu hình được chạy trên năm seed và báo cáo
dưới dạng trung bình kèm độ lệch chuẩn.

%%===========================================================%%
\section{Kết quả}\label{sec:results}

\subsection{So sánh tại mức dịch chuyển đại diện}
Tại mức dịch chuyển được lựa chọn cho đánh giá chính, coverage của phương pháp hiệu chỉnh tĩnh giảm từ
$90.5\%$ trên tập hiệu chỉnh xuống còn $81.0\%$ trên tập kiểm tra, cho thấy sự xuất hiện rõ rệt của dịch
chuyển phân phối. Trong khi đó, độ chệch của backbone gần như bằng không và giá trị $\sigma$-gain dương
($+0.44$), xác nhận rằng thang phương sai Poisson--Binomial vẫn tương quan với sai số thực. Kết quả so sánh
được trình bày trong Bảng~\ref{tab:main}.

\begin{table}[h]
\caption{So sánh joint $K{=}2$ tại một mức dịch chuyển cố định. Coverage và Winkler báo cáo dưới dạng trung
bình $\pm$ độ lệch chuẩn trên $5$ seed, AvgW là giá trị trung bình; Winkler càng thấp càng tốt.}\label{tab:main}
\begin{tabular}{@{}lccc@{}}
\toprule
Phương pháp & Coverage\,(\%) & AvgW & Winkler \\
\midrule
Static split-CP & $81.0 \pm 0.0$ & 0.03 & $2.84 \pm 0.00$ \\
ACI~\cite{gibbs2021adaptive} & $89.7 \pm 0.1$ & 6.29 & $14.08 \pm 2.66$ \\
NexCP~\cite{barber2023conformal} & $88.0 \pm 0.7$ & 0.21 & $2.23 \pm 0.06$ \\
FACI~\cite{gibbs2024conformal} & $88.6 \pm 0.7$ & 0.26 & $2.19 \pm 0.06$ \\
SAOCP~\cite{bhatnagar2023improved} & $99.3 \pm 0.2$ & 11.12 & $22.34 \pm 0.54$ \\
COP~\cite{hu2026conformal} & $85.1 \pm 0.5$ & 0.10 & $2.52 \pm 0.04$ \\
Rolling-Origin~\cite{halkiewicz2026rolling} & $88.5 \pm 0.3$ & 0.24 & $2.18 \pm 0.04$ \\
PB-JCI Fixed & $85.7 \pm 0.9$ & 0.18 & $2.32 \pm 0.02$ \\
\textbf{Adaptive PB-JCI} & $\mathbf{91.2 \pm 0.3}$ & 0.47 & $\mathbf{2.10 \pm 0.09}$ \\
\botrule
\end{tabular}
\end{table}

Bảng~\ref{tab:main} cho thấy Adaptive PB-JCI là phương pháp duy nhất đồng thời đạt coverage không thấp hơn
mức mục tiêu ($91.2\%$) và điểm Winkler thấp nhất ($2.10$). Các phương pháp còn lại hoặc under-cover, hoặc
đạt coverage bằng các khoảng dự đoán rộng hơn đáng kể (SAOCP, ACI).

\subsection{Ổn định dưới cường độ dịch chuyển tăng dần}
Bảng~\ref{tab:depth} trình bày joint coverage khi cường độ dịch chuyển tăng dần. Trong thiết lập này,
coverage của phương pháp hiệu chỉnh tĩnh giảm từ $85.5\%$ xuống còn $9.0\%$, phản ánh sự suy giảm nghiêm
trọng của hiệu chỉnh cố định khi khoảng cách giữa dữ liệu hiệu chỉnh và dữ liệu triển khai ngày càng lớn.

\begin{table}[h]
\caption{Joint coverage (\%) theo cường độ dịch chuyển. Adaptive PB-JCI luôn không thấp hơn mức mục tiêu
trong khi các baseline suy giảm hoặc over-cover.}\label{tab:depth}
\begin{tabular}{@{}lccccccccc@{}}
\toprule
Dịch chuyển & Static & ACI & NexCP & FACI & SAOCP & COP & Roll. & PB-Fix & \textbf{Adapt.} \\
\midrule
Nhẹ        & 85.5 & 89.8 & 88.7 & 88.8 & 98.3 & 87.3 & 89.2 & 87.9 & \textbf{91.4} \\
Vừa        & 81.0 & 89.7 & 88.0 & 88.6 & 99.3 & 85.1 & 88.5 & 85.7 & \textbf{91.2} \\
Mạnh       & 53.1 & 89.6 & 87.1 & 88.6 & 99.3 & 82.9 & 87.0 & 83.3 & \textbf{91.0} \\
Rất mạnh   & 9.0  & 89.7 & 86.2 & 89.1 & 99.5 & 80.3 & 85.8 & 81.9 & \textbf{90.5} \\
\botrule
\end{tabular}
\end{table}

Ngược lại, Adaptive PB-JCI duy trì coverage ổn định quanh mức mục tiêu ($90.5$--$91.4\%$) trên toàn bộ dải
khảo sát. Xu hướng này được thể hiện trực quan trong Hình~\ref{fig:depth}.

\begin{figure}[h]
\centering
\includegraphics[width=0.72\textwidth]{fig_pub_depth}
\caption{Joint coverage theo cường độ dịch chuyển (\texttt{DEPTH}). Adaptive PB-JCI (xanh) bám sát mục tiêu
$90\%$ (đường ngang nét đứt); SAOCP (cam) over-cover; hiệu chỉnh tĩnh (xám đậm, nét đứt) sụp đổ xuống $9\%$;
các baseline còn lại (xám nhạt) lệch dần khỏi mục tiêu.}\label{fig:depth}
\end{figure}

\subsection{Động học khôi phục coverage}
Hình~\ref{fig:recovery} minh họa coverage cuộn trên một dòng dữ liệu chuyển đột ngột từ khảo sát sâu sang
khảo sát nông.

\begin{figure}[h]
\centering
\includegraphics[width=0.9\textwidth]{fig_pub_recovery}
\caption{Coverage cuộn sau một dịch chuyển đột ngột sâu$\to$nông (đường đứng: thời điểm shift; nền nhạt: giai
đoạn khảo sát nông). Adaptive PB-JCI (xanh đậm) khôi phục về mức mục tiêu $90\%$ trong khoảng $30$ bước; hiệu
chỉnh tĩnh (xám đậm, nét đứt) suy giảm kéo dài và SAOCP (cam) over-cover, các baseline còn lại vẽ mờ màu xám
nhạt.}\label{fig:recovery}
\end{figure}

Sau điểm chuyển, các phương pháp ngưỡng cố định suy giảm coverage kéo dài còn SAOCP over-cover, trong khi
Adaptive PB-JCI khôi phục coverage về gần mức mục tiêu trong khoảng $30$ bước và duy trì ổn định sau đó.

%%===========================================================%%
\section{Thảo luận}\label{sec:discussion}

Kết quả cho thấy Adaptive PB-JCI phát huy hiệu quả khi backbone có độ chệch thấp, các xác suất đầu ra phản
ánh đúng mức độ bất định (thể hiện qua $\sigma$-gain dương) và dịch chuyển phân phối xuất hiện theo hướng làm
suy giảm chất lượng đếm. Trong điều kiện đó, thang phương sai Poisson--Binomial và cơ chế hiệu chỉnh trực
tuyến bổ trợ cho nhau để duy trì coverage mục tiêu. Các kết quả vẫn giữ nguyên về mặt định tính khi $\kappa$
được thay đổi trong khoảng $[1.0, 3.0]$.

Nghiên cứu này tập trung vào hiệu chỉnh số đếm trên một danh mục SDSS cố định. Do nhãn \texttt{probPSF} của
SDSS là nhị phân và danh mục chỉ chứa các nguồn đã phát hiện, cả xác suất phát hiện lẫn độ mờ phân loại dưới
dịch chuyển đều là thành phần được mô hình hóa (neo trên SNR thật), nên benchmark mang tính bán tổng hợp. Một
đánh giá hoàn toàn quan trắc, chẳng hạn dùng xác suất phân loại liên tục hoặc so sánh giữa nhiều khảo sát
khác nhau, là hướng mở rộng tự nhiên trong tương lai.

%%===========================================================%%
\section{Kết luận}\label{sec:conclusion}

Nghiên cứu này khảo sát khả năng chuyển giao Adaptive PB-JCI Online sang bài toán đếm nguồn thiên văn dưới
dịch chuyển phân phối. Trên dữ liệu SDSS DR17, phương pháp là lựa chọn duy nhất duy trì joint coverage ở mức
mục tiêu đồng thời đạt điểm Winkler thấp nhất trong số các phương pháp được đánh giá. Kết quả cho thấy các kỹ
thuật hiệu chỉnh đếm có bảo đảm coverage có thể được chuyển giao hiệu quả từ khoa học sự sống sang vật lý
thiên văn, và có tiềm năng áp dụng cho các bài toán đếm khác dưới dịch chuyển phân phối.

%%===========================================================%%
\backmatter

\bmhead{Declarations}
Đây là công trình gốc, chưa được công bố và không đồng thời được nộp ở nơi khác. Về đóng góp của các tác giả,
T.T.H.~Dinh thực hiện cài đặt phương pháp, tiến hành thí nghiệm, phân tích và viết bản thảo, còn V.H.~Duong
phụ trách ý tưởng, hướng dẫn và chỉnh sửa. Về việc sử dụng công cụ trí tuệ nhân tạo, nhóm tác giả có dùng các
công cụ hỗ trợ trong soạn thảo văn bản và viết mã; toàn bộ thiết kế thí nghiệm, dữ liệu, kết quả và kết luận
do các tác giả thực hiện, kiểm chứng và chịu trách nhiệm, phù hợp với yêu cầu công khai của ICEBA2026. Dữ
liệu SDSS DR17 được cung cấp công khai qua CasJobs và astroquery. Các tác giả tuyên bố không có xung đột lợi
ích.

%%===========================================================%%
\bibliography{references}

\end{document}
